\section{Related Work}

Automated-vehicle evaluation literature has emphasized for years that scenario-based simulation is not just a convenience layer, but a core methodological instrument for establishing trustworthy performance and safety claims \cite{riedmaier2020scenario}. That literature typically focuses on coverage, scenario selection, and validation logic at the vehicle-system level. Our work is aligned with that concern, but addresses a newer source of evaluation fragility: the runtime instability introduced when the driving policy itself is a locally deployed large language model.

Recent work has established that large language models can act as high-level decision components for autonomous driving, but the literature is still split between papers that emphasize behavioral capability and papers that emphasize interpretability or system design. Early work such as \emph{Drive Like a Human} argued that language models can serve as driving-policy interfaces rather than pure text processors \cite{fu2023drivelikehuman}. Follow-up systems including DiLu, LanguageMPC, and DriveGPT4 broadened that direction by treating language models as planners, decision makers, or interpretable end-to-end controllers in simulated driving settings \cite{wen2023dilu,sha2023languagempc,xu2023drivegpt4}. A recent survey, \emph{LLM4Drive}, makes clear that this area is expanding quickly across planning, perception-language interaction, and embodied control \cite{yang2023llm4drive}. Our paper is closest to this line of work, but differs in emphasis: we do not primarily ask whether an LLM driving agent can complete a scenario. We ask when the local evaluation regime remains valid enough for those comparisons to be trusted.

That distinction also changes how we use the fine-tuning literature. Parameter-efficient adaptation methods such as LoRA and QLoRA make it feasible to specialize local language models without full retraining \cite{hu2021lora,dettmers2023qlora}. Systems work on local inference, including PowerInfer, further shows that practical deployment quality depends heavily on hardware-aware runtime behavior rather than model capability in isolation \cite{song2023powerinfer}. Broader evaluation work such as HELM argues that meaningful model comparison requires multiple dimensions of assessment rather than single-score reporting \cite{liang2022helm}. DiLu-Ollama sits at the intersection of these threads. It studies fine-tuned local driving agents, but evaluates them through a benchmark that preserves invalid runs, timeout collapse, and latency-induced comparability failures.
